% ==================================================
%  Preámbulo — Reporte de Consultoría: Metalentes
% ==================================================

% --- Codificación y lenguaje ---
\usepackage[T1]{fontenc}
\usepackage[utf8]{inputenc}
\usepackage[spanish]{babel}
\decimalpoint
% Rotula los flotantes de tabla como "Tabla" (no "Cuadro"), coherente con "Figura"
\addto\captionsspanish{\renewcommand{\tablename}{Tabla}}
\usepackage{lmodern}

% --- Matemáticas ---
\usepackage{amsmath, amssymb, mathtools}

% --- Diseño de página ---
\usepackage[left=30mm, right=30mm, top=30mm, bottom=30mm]{geometry}
\usepackage{parskip}
\setlength{\parskip}{6pt}
\setlength{\emergencystretch}{3em}

% --- Interlineado (espaciado entre líneas un poco mayor) ---
\usepackage{setspace}
\setstretch{1.3}

% --- Elementos visuales ---
\usepackage{booktabs, array, multicol}
\usepackage{xcolor}
\usepackage{graphicx, float}
\usepackage{fancyhdr}

% --- Leyendas de tablas y figuras (más pequeñas que el cuerpo, etiqueta en negrita) ---
\usepackage{caption}
\captionsetup{font=footnotesize, labelfont=bf}
\usepackage[hidelinks]{hyperref}
\usepackage{tikz}          % Necesario para la portada
\usetikzlibrary{babel}     % Evita conflictos entre TikZ y babel (español)
\usepackage{standalone}    % Permite \input de archivos .tikz standalone
\usepackage{subcaption}    % Subfiguras lado a lado

% --- Código fuente --
\usepackage{listings}
\lstset{
  basicstyle=\ttfamily\footnotesize,
  breaklines=true,
  frame=single,
  language=Python,
  numbers=left,
  numberstyle=\tiny\color{gray},
  keywordstyle=\color{blue},
  commentstyle=\color{green!50!black},
  stringstyle=\color{red!70!black},
}

% --- Cajas de énfasis ---
\usepackage[skins,breakable]{tcolorbox}

% --- Tablas largas (que se parten entre páginas) y rutas de archivo partibles ---
\usepackage{longtable}
\usepackage{url}
\urlstyle{tt}
\makeatletter
\g@addto@macro\UrlBreaks{\do\/\do\_\do\.\do\-}
\makeatother
\newcommand{\fpath}[1]{\nolinkurl{#1}}  % ruta monoespaciada que rompe en / _ . -

% --- Bibliografía ---
\usepackage{csquotes}
\usepackage[backend=biber, style=numeric, sorting=none]{biblatex}
\addbibresource{Bibliografia/referencias.bib}
% Fija el tamaño de letra de la bibliografía al del cuerpo del documento
% (el cuerpo hereda \footnotesize de la portada; sin esto, las entradas
% se ven mayores que el texto de los capítulos).
\renewcommand*{\bibfont}{\normalfont\footnotesize}
